\section{Contributions}
\label{sec:introduction_contributions}

This thesis makes three main contributions.

First, it presents the design and implementation of a VR indoor rowing prototype that supports two distinct representation strategies for the same physical rowing activity. The prototype distinguishes between an ergometer-congruent mode and a boat-centric mode. This distinction makes it possible to examine a design trade-off that is often implicit in XR sport systems: whether virtual representation should prioritize the user's actual input apparatus and body movement, or the visual form of the sport being simulated.

Second, the thesis contributes an embodiment-oriented comparison of these two representation strategies. Rather than treating immersion, realism, or motivation as broad and undifferentiated goals, the thesis focuses on how movement mapping may affect specific components of embodied experience. In particular, it considers sense of agency, body ownership, self-location, and overall sense of embodiment as relevant constructs for evaluating rowing representation in VR.

Third, the thesis provides design implications for researchers and developers of XR sport and exercise systems in which physical equipment and virtual sport depiction do not fully coincide. The rowing ergometer is a clear example of this situation, but the underlying problem also applies to other stationary training systems where real input devices are used to represent more complex sport actions. The thesis therefore contributes to a broader discussion about how embodied XR systems should balance sensorimotor congruence, visual plausibility, sport-specific realism, and user preference.

The intended contribution is not a general claim that one representation strategy is universally superior. Instead, the thesis aims to clarify what is gained and what may be lost when a VR sport system chooses between direct physical congruence and realistic sport depiction. Both strategies solve part of the design problem while introducing another. A congruent representation may feel controllable and bodily grounded, while a sport-realistic representation may feel more meaningful or recognizable as rowing. Understanding this trade-off can help future systems make more deliberate representation choices.

In practical terms, the thesis also documents the technical and methodological considerations involved in building such a comparison. These include alignment between physical and virtual rowing spaces, avatar motion mapping, control of visual confounds, and the selection of measures for evaluating embodiment and experience. These details are part of the contribution because the comparison depends not only on the final visual result, but on the system's ability to isolate the relevant design difference.
